% Overleaf: izabrati XeLaTeX. Sve skice su u ovom fajlu.
\documentclass[10pt,a4paper]{article}
\usepackage[margin=20mm,headheight=14pt]{geometry}
\usepackage{fontspec}
\setmainfont{DejaVu Serif}[ItalicFont={DejaVu Serif},ItalicFeatures={FakeSlant=0.2}]
\setsansfont{DejaVu Sans}
\setmonofont{DejaVu Sans Mono}
\usepackage{amsmath,amssymb,booktabs,array,microtype}
\usepackage{tikz,pgfplots}
\usetikzlibrary{arrows.meta,positioning,calc,fit}
\pgfplotsset{compat=1.18}
\usepackage{xurl,hyperref,fancyhdr}
\usepackage{polyglossia}
\setdefaultlanguage[Script=Cyrillic]{serbian}
\setlength{\emergencystretch}{3em}
\hypersetup{colorlinks=true,urlcolor=blue!50!black,linkcolor=blue!50!black}
\pagestyle{fancy}\fancyhf{}
\fancyhead[L]{\small Рачунарски вид — припрема за усмени испит}
\fancyfoot[C]{\thepage}
\setlength{\parindent}{0pt}\setlength{\parskip}{4pt}
\setlength{\abovedisplayskip}{5pt}\setlength{\belowdisplayskip}{5pt}
\newcommand{\en}[1]{\textit{#1}}
\newcommand{\Sze}{Szeliski, \en{Computer Vision: Algorithms and Applications}, 2. издање}
\newcounter{pitanje}
\newenvironment{odgovor}[1]{\stepcounter{pitanje}\begin{minipage}[t][0.455\textheight][t]{\textwidth}\vspace{0pt}\textbf{\thepitanje. #1}\par\smallskip}{\end{minipage}\par\ifodd\value{pitanje}\vspace{0.025\textheight}\hrule\vspace{0.015\textheight}\else\newpage\fi}
\newcommand{\izvor}[1]{\par\vfill{\footnotesize\textbf{За читање:} #1}\par}
\tikzset{box/.style={draw,rounded corners=2pt,fill=blue!5,align=center,font=\small,inner sep=4pt},arr/.style={-{Latex[length=2mm]},thick}}
\newcommand{\lanac}[1]{\begin{center}\small #1\end{center}}

\newcommand{\projsketch}{\begin{center}\begin{tikzpicture}[scale=.65]
\draw[arr](-.5,0)--(9,0) node[right]{$Z$};\draw[arr](0,-.2)--(0,3.7) node[above]{$X$};
\draw[thick,blue](2,-.6)--(2,3.2) node[above]{раван $Z=f$};
\fill (0,0) circle(2pt) node[below]{$C$};\fill (7,3) circle(2pt) node[above]{$P=(X,Z)$};
\draw[arr](0,0)--(7,3);\fill[blue](2,.857) circle(3pt) node[right]{$x=fX/Z$};
\draw[dashed](7,3)--(7,0);\draw[<->](0,-.65)--node[below]{$f$}(2,-.65);
\end{tikzpicture}\end{center}}
\newcommand{\episketch}{\begin{center}\begin{tikzpicture}[scale=.65]
\fill[blue!5] (0,0)--(8,0)--(4,3.4)--cycle;
\draw[fill=gray!10] (2,-.6)--(2.6,-.9)--(2.6,2.5)--(2,2.8)--cycle;
\draw[fill=gray!10] (6,-.6)--(6.6,-.9)--(6.6,2.5)--(6,2.8)--cycle;
\draw[thick](0,0)--(4,3.4)--(8,0);\draw[dashed](0,0)--(8,0);
\draw[very thick,blue](2,-.5)--(2,2.6) node[above]{$\ell$};
\draw[very thick,red](6,-.5)--(6,2.6) node[above]{$\ell'$};
\foreach \x/\y/\label in {0/0/C_1,8/0/C_2,4/3.4/X,2/0/e,6/0/e',2/1.7/x,6/1.7/x'}{\fill(\x,\y) circle(2pt) node[above right]{$\label$};}
\node at (4,-.65){основица $C_1C_2$};\node at (4,2.5){$\Pi$};
\end{tikzpicture}\end{center}}
\begin{document}
\begin{center}
{\LARGE\bfseries Рачунарски вид}\\[6pt]
{\Large Скрипта за усмени испит}\\[6pt]
{\large Одговори на 50 испитних питања}
\end{center}
\medskip
Скрипта прати редослед питања из документа \texttt{CV\_Ispitna\_pitanja.pdf}. Одговори садрже суштину за усмено излагање, најважније формуле и скице. Формулације питања су правописно уједначене. Свака страница са питањима садржи два одговора; један одговор, заједно са питањем и референцом, заузима највише приближно пола странице.

\textbf{Главна литература [S].} Richard Szeliski, \en{Computer Vision: Algorithms and Applications}, \textbf{2. издање}, Springer, 2022. Испод сваког одговора наведени су конкретни одељци за читање из ове књиге. Њихови називи задржани су на енглеском ради лакшег проналажења. За специфичну улогу \en{Mean Shift}-а у Kinect систему додат је изворни рад [A1], наведен на крају документа.

\textbf{Конвенције.} $I$ је интензитет слике; $x,y$ су просторне координате пиксела, а $u,v$ координате пројекције у геометрији камере. $X,Y,Z$ означавају координате 3D тачке у систему камере. $K$ је матрица унутрашњих параметара. Користи се виртуелна раван слике \textbf{испред} центра камере и позитивно $Z$ ка сцени, па у пројекцији нема знака минус. Симбол $\sim$ значи једнакост до ненултог фактора у хомогеним координатама.

Код конволуције и \en{pooling}-а, ако није другачије речено, попуна је нула, а корак је 1. У морфолошким скицама плаво означава објекат, а бело позадину. Енглески технички термини у одговорима су у курзиву.

\textbf{Коришћење.} За сваки одговор прво објаснити идеју, затим написати формулу или нацртати скицу. Код питања чији одговор зависи од модела или претпоставки, те претпоставке су експлицитно наведене. Наведени извори служе за детаљније разумевање, а скрипта је сажета припрема.

% \textbf{Overleaf.} Копирати овај фајл у \texttt{main.tex} и изабрати \texttt{XeLaTeX}. Све скице су унутар документа; нису потребне спољне слике ни \texttt{.bib} фајл.
\newpage
\begin{odgovor}{Како детектовати ивице на слици са шумом коришћењем најмање два оператора?}
Најпре применимо Гаусов оператор изглађивања да смањимо шум (енг. \textit{Gaussian blur, smoothing}), а затим оператор извода, нпр. \en{Sobel}, да издвојимо промене интензитета:
\[
I_\sigma=G_\sigma*I,\qquad I_x=S_x*I_\sigma,\quad I_y=S_y*I_\sigma,\qquad
M=\sqrt{I_x^2+I_y^2}.
\]
\[
S_x=\begin{bmatrix}-1&0&1\\-2&0&2\\-1&0&1\end{bmatrix},\qquad S_y=S_x^\top.
\]
Велика вредност у $M$ означава кандидата за ивицу; праг (енг. \textit{threshold}) одваја јаке промене од слабих. Знак извода зависи од избора корелације или конволуције, али то не утиче на његову магнитуду.

Алтернатива је \en{Laplacian of Gaussian} — LoG: изглађивање па Лапласијан, уз тражење промене знака, односно \en{енг. zero crossings} --- нула функције. Операције се могу спојити јер $\nabla^2(G_\sigma*I)=(\nabla^2G_\sigma)*I$. Сам извод без претходног изглађивања појачава шум. Веће $\sigma$ боље потискује шум, али може изгубити ситне ивице.
\izvor{Прочитати §3.2.2, \en{Examples of linear filtering}, и §7.2.1, \en{Edge detection}, из [S].}
\end{odgovor}

\begin{odgovor}{Која је разлика између Гаусовог шума и salt-and-pepper шума и који филтри их уклањају?}
Гаусов шум се често моделира као адитиван:
\[
I_{\mathrm{noisy}}(x,y)=I(x,y)+n(x,y),\qquad n\sim\mathcal N(0,\sigma_n^2).
\]
Многи пиксели добијају мале случајне промене, са Гаусовом расподелом амплитуда. Гаусово или просечно изглађивање смањује такав шум, али замућује ивице. \en{Bilateral filter} комбинује просторну близину и сличност интензитета, па боље чува ивице.

Шум \en{salt-and-pepper} је импулсни: поједини пиксели замењују се минималном или максималном вредношћу, нпр. 0 и 255. \textbf{Медијански филтар} замењује пиксел медијаном његовог суседства, па је отпоран на такве екстреме. На пример, медијана $[98,100,101,102,255]$ је 101, док просек 131.2 остаје снажно повучен импулсом. Не постоји један најбољи филтар за све врсте шума: избор зависи од модела шума и детаља које треба сачувати.
\izvor{Прочитати §3.2.2, §3.3.1, \en{Non-linear filtering}, и §3.3.2, \en{Bilateral filtering}, из [S].}
\end{odgovor}

\begin{odgovor}{Који филтар се користи за sharpening и да ли је битно користи ли се конволуција или корелација?}
За пооштравање можемо додати високофреквентни део слике:
\[
I_{\mathrm{sharp}}=I+\alpha(I-G_\sigma*I),\qquad \alpha>0.
\]
Ово је \en{unsharp masking}: од оригинала одузимамо замућену слику, па разлику додајемо оригиналу. Лапласијански филтар такође појачава промене; за
\[
L=\begin{bmatrix}0&1&0\\1&-4&1\\0&1&0\end{bmatrix}
\]
можемо користити $I_{\mathrm{sharp}}=I-L*I$, односно језгро
\[
H=\begin{bmatrix}0&-1&0\\-1&5&-1\\0&-1&0\end{bmatrix}.
\]
Конволуција обрће језгро за $180^\circ$, а корелација га не обрће. Наведено језгро је централно симетрично, па даје исти резултат у оба случаја, уз исте граничне услове. За произвољно несиметрично језгро резултати се разликују. Пооштравање појачава и шум и може створити светле или тамне ореоле око ивица.
\izvor{Прочитати §3.2, \en{Linear filtering}, §3.2.2 и §3.4.2, \en{Application: Sharpening, blur, and noise removal}, из [S].}
\end{odgovor}

\begin{odgovor}{Како се морфолошко отварање и затварање дефинишу преко дилатације и ерозије? Графички илустровати.}
За бинарни објекат $A$ и структурни елемент $B$:
\[
A\circ B=(A\ominus B)\oplus B,\qquad A\bullet B=(A\oplus B)\ominus B.
\]
Ерозија $\ominus$ задржава позиције у којима цео структурни елемент стаје у објекат, а дилатација $\oplus$ проширује објекат његовим померањем. \textbf{Отварање} прво еродира па дилатира: уклања ситне објекте и танке израштаје. \textbf{Затварање} прво дилатира па еродира: попуњава мале рупе и прекиде и може спојити блиске компоненте.
\begin{center}\begin{tikzpicture}[scale=.75]
\node at (0,2){Отварање};\node at (4,2){Почетно};\node at (8,2){Затварање};
\foreach \x in {0,4,8}{\draw[fill=blue!35] (\x-1,0) rectangle (\x+1,1.3);}
\draw[fill=white](-.25,.45) rectangle(.25,.85);\fill[blue!35](5.4,.8) circle(.1);
\draw[fill=white](3.75,.45) rectangle(4.25,.85);\fill[blue!35](9.4,.8) circle(.1);
\draw[arr](2.7,.5)--(1.7,.5);\draw[arr](5.7,.5)--(6.7,.5);
\node[font=\scriptsize] at (0,-.4){ситна компонента нестаје};
\node[font=\scriptsize] at (8,-.4){мала рупа нестаје};
\end{tikzpicture}\end{center}
Скице су шематске и исход зависи од величине и облика $B$. Редослед операција је битан; отварање и затварање нису међусобни инверзи.
\izvor{Прочитати §3.3.3, \en{Binary image processing}, из [S], део о математичкој морфологији.}
\end{odgovor}

\begin{odgovor}{Какав ефекат даје већи структурни елемент код затварања у односу на мањи? Графички илустровати.}
Затварање са већим структурним елементом може попунити веће рупе и затворити шире прекиде, па лакше спаја блиске компоненте. Мањи елемент мења само ситне детаље. Цена већег елемента је губитак реалних отвора и раздвајања између објеката.
\begin{center}\begin{tikzpicture}[scale=.8]
\foreach \y/\lab in {2/Улаз,1/Мали $B$,0/Већи $B$}{\node[anchor=east] at (-.5,\y+.25){\lab};}
\foreach \y in {1,2}{\draw[fill=blue!35](0,\y) rectangle(2,\y+.5);\draw[fill=blue!35](2.8,\y) rectangle(4.8,\y+.5);}
\draw[fill=blue!35](0,0) rectangle(4.8,.5);
\draw[<->](2,2.7)--node[above,font=\small]{прекид}(2.8,2.7);
\node[font=\small,anchor=west] at (5.2,1.25){прекид остаје};
\node[font=\small,anchor=west] at (5.2,.25){прекид се затвара};
\end{tikzpicture}\end{center}
Скица приказује хоризонтални структурни елемент: кратак не премошћује прекид, а дужи га премошћује. Важан је и \textbf{облик}: дугачак хоризонтални елемент не мора затворити вертикални прекид исте ширине. За уобичајене угнежђене елементе истог облика, већа размера значи деловање на крупније структуре; није реч само о броју пиксела у елементу.
\izvor{Прочитати §3.3.3, \en{Binary image processing}, из [S], операције дилатације, ерозије и затварања.}
\end{odgovor}

\begin{odgovor}{Како изгледа Фуријеов спектар слике са хоризонталним синусоидалним променама интензитета? Скицирати и објаснити.}
Ако се интензитет мења \textbf{дуж хоризонталне координате}:
\[
I(x,y)=a+b\cos(2\pi f_0x+\phi),
\]
спектар има DC компоненту у $(0,0)$ ако је $a\ne0$ и два симетрична врха на $(f_0,0)$ и $(-f_0,0)$. Они представљају позитивну и негативну фреквенцију реалног синусоидалног сигнала.
\begin{center}\begin{tikzpicture}[scale=.8]
\draw[arr](-3,0)--(3.5,0) node[right]{$f_x$};
\draw[arr](0,-1)--(0,1.5) node[above]{$f_y$};
\fill[blue](-2,0) circle(4pt) node[below]{$-f_0$};
\fill[blue](2,0) circle(4pt) node[below]{$f_0$};
\fill(0,0) circle(3pt) node[below right]{DC};
\end{tikzpicture}\end{center}
Фаза $\phi$ мења комплексне фазе врхова, а не њихове положаје. Ако „хоризонталне промене“ означавају \textbf{хоризонталне пруге}, онда се интензитет мења по $y$, па су врхови на вертикалној оси $(0,\pm f_0)$. Код DFT-а најчистији врхови настају када фреквенција одговара дискретном фреквенцијском бину; иначе се јавља \en{spectral leakage}.
\izvor{Прочитати §3.4, \en{Fourier transforms}, и §3.4.1, \en{Two-dimensional Fourier transforms}, из [S].}
\end{odgovor}

\begin{odgovor}{Зашто конволуција слике Гаусовим филтром смањује високе фреквенције?}
Конволуција у просторном домену одговара множењу спектра. За нормализовани Гаусов филтар стандардне девијације $\sigma$:
\[
G_\sigma(x,y)=\frac1{2\pi\sigma^2}e^{-(x^2+y^2)/(2\sigma^2)},\qquad
\widehat G_\sigma(f_x,f_y)=e^{-2\pi^2\sigma^2(f_x^2+f_y^2)}.
\]
Стога је
\[
\widehat I_{\mathrm{out}}=\widehat G_\sigma\,\widehat I.
\]
Фактор је близу 1 за ниске фреквенције и опада ка нули за високе: филтар је \en{low-pass}. Брзе просторне промене, као што су ситан шум, танке структуре и оштре ивице, потискују се. У просторном домену то је пондерисано усредњавање суседних пиксела.

Веће $\sigma$ даје шири филтар у простору и ужу пропусну област у спектру, па јаче замућење. Гаусово филтрирање не уклања све високе фреквенције идеалним одсецањем, већ их постепено пригушује.
\izvor{Прочитати §3.2.1, \en{Separable filtering}, §3.2.2 и §3.4, \en{Fourier transforms}, из [S].}
\end{odgovor}

\begin{odgovor}{Како се манифестује aliasing, шта га узрокује и како се избегава? Графички илустровати.}
\en{Aliasing} настаје када узорковање није довољно густо за фреквенције сигнала. Различити континуални сигнали тада дају исте узорке. Последице су лажни нискофреквентни обрасци, \en{moiré}, назубљене ивице и погрешне текстуре.
\begin{center}\begin{tikzpicture}[x=1.8cm,y=.9cm]
\draw[arr](-.12,0)--(5.25,0) node[right]{$x$};
\draw[arr](0,-1.2)--(0,1.3);
\foreach \x in {1,2,3,4,5}{\draw(\x,.05)--(\x,-.05) node[below,font=\scriptsize]{$\x$};}
\node[above left,font=\scriptsize] at(0,1){1};\node[below left,font=\scriptsize] at(0,-1){$-1$};
\draw[blue,thick,domain=0:5,samples=220] plot(\x,{cos(360*.8*\x)});
\draw[red,dashed,thick,domain=0:5,samples=160] plot(\x,{cos(360*.2*\x)});
\foreach \x/\y in {0/1,1/.309,2/-.809,3/-.809,4/.309,5/1}{\fill(\x,\y) circle(1.8pt);}
\end{tikzpicture}\end{center}
Плави сигнал фреквенције 0.8 и црвени 0.2 дају исте узорке на целобројним позицијама. За сигнал ограничен на $f_{\max}$ потребно је узорковати брже од $2f_{\max}$, уз одговарајуће ограничење спектра.

Пре смањења резолуције примењује се \en{anti-aliasing low-pass filter}, нпр. Гаусов, прилагођен новом кораку узорковања. Само прескакање пиксела није довољно; већ настали \en{aliasing} не може се уопштено уклонити накнадним замућивањем.
\izvor{Прочитати §2.3.1, \en{Sampling and aliasing}, и §3.5.2, \en{Decimation}, из [S].}
\end{odgovor}

\begin{odgovor}{Описати детекцију ивица Canny детектором.}
\en{Canny} спаја изглађивање, процену градијента, стањивање и повезивање ивица:
\lanac{Гаусово изглађивање $\to$ градијент $\to$ \en{non-maximum suppression} $\to$ два прага $\to$ \en{hysteresis}}
Прво се смањује шум и рачунају величина градијента $M=\sqrt{I_x^2+I_y^2}$ и смер $\theta=\operatorname{atan2}(I_y,I_x)$. \en{Non-maximum suppression} задржава само локалне максимуме величине \textbf{дуж смера градијента}, који је нормалан на ивицу. Тако се широк одзив претвара у танку ивицу.

Са праговима $T_{\mathrm{low}}<T_{\mathrm{high}}$, пиксели изнад високог прага су јаке ивице, испод ниског се одбацују, а средњи су слаби кандидати. \en{Hysteresis} задржава слабе кандидате само ако су повезани са јаким ивицама, преко ланца кандидата. Зато је резултат мање испрекидан него код једног прага. Величина изглађивања и прагови контролишу компромис између шума и пропуштених ивица.
\izvor{Прочитати §7.2.1, \en{Edge detection}, из [S], део о Canny детектору, локалним максимумима и хистерези.}
\end{odgovor}

\begin{odgovor}{Који филтар, односно функција, може издвојити ивице на слици?}
Ивица је место нагле промене интензитета, па је издвајамо \textbf{изводом}. Први извод се процењује филтрима \en{Sobel}, \en{Prewitt} или изводима Гаусове функције. За изглађену слику:
\[
I_x=(\partial_xG_\sigma)*I,\qquad I_y=(\partial_yG_\sigma)*I,\qquad M=\sqrt{I_x^2+I_y^2}.
\]
Ивице су кандидати на локалним максимумима $M$. Извод Гаусове функције истовремено уводи скалу изглађивања и мери промену.

Други приступ користи Лапласијан $\nabla^2I=I_{xx}+I_{yy}$, обично после изглађивања, односно \en{LoG}. Ивице се тада траже на променама знака одзива, уз проверу његове јачине. \en{Difference of Gaussians} — DoG — за блиске скале приближно представља скалирани LoG. Сам Гаусов филтар изглађује слику; његов извод или разлика скала издвајају промене.
\izvor{Прочитати §3.2.3, \en{Band-pass and steerable filters}, и §7.2.1, \en{Edge detection}, из [S].}
\end{odgovor}

\begin{odgovor}{Зашто се Hough трансформација заснива на гласању за параметарске облике и зашто то даје резултате?}
\en{Hough transform} претвара откривање облика у тражење максимума у параметарском простору. За праву је погодан облик
\[
\rho=x\cos\theta+y\sin\theta.
\]
Сваки ивични пиксел $(x,y)$ гласа за све парове $(\rho,\theta)$ који описују праву кроз њега. У параметарском простору његови гласови леже на синусоиди. Пиксели са исте праве имају криве које се секу код заједничких параметара, па настаје јак максимум акумулатора.
\lanac{Ивични пиксели $\to$ криве у $(\rho,\theta)$ $\to$ врх акумулатора $\to$ детектована права}
Техника омогућава да удаљени или испрекидани делови једног облика заједнички допринесу детекцији. Насумичан шум обично не гласа доследно за исти модел. Дискретизација одређује прецизност и меморију; јаки случајни обрасци могу ипак створити лажне врхове. За кругове се гласа за центар и полупречник, па расте димензија претраге.
\izvor{Прочитати §7.4.2, \en{Hough transforms}, из [S].}
\end{odgovor}

\begin{odgovor}{За налажење билијарских лоптица, да ли је прикладнији Hough или RANSAC? Образложити.}
За више јасно видљивих, приближно кружних лоптица са сличним величинама, \textbf{кружна Hough трансформација} је природан почетни избор:
\[
(x-a)^2+(y-b)^2=r^2.
\]
Ивични пиксели гласају за $(a,b,r)$, а више врхова акумулатора даје више лоптица. Познат распон полупречника и смер градијента могу знатно ограничити претрагу. Пре тога могу се издвојити области чија боја одступа од боје стола.

\en{RANSAC} је погодан ако већ имамо кандидате за ивичне тачке са много погрешних тачака или делимично заклоњене кругове: три неколинеарне тачке одређују круг, а остале проверавају консензус. За више лоптица треба поновити издвајање модела или користити више-моделски поступак.

Избор није универзалан. Јака перспектива може кружни обрис претворити у елипсу, па модел круга није адекватан; потребна је ректификација одговарајуће равни или сложенији модел обриса.
\izvor{Прочитати §7.4.2, \en{Hough transforms}, и §8.1.4, \en{Robust least squares and RANSAC}, из [S].}
\end{odgovor}

\begin{odgovor}{Шта значи да дескриптори кључних тачака треба да буду invariant и distinctive? Зашто је инваријантност неопходна?}
\textbf{Инваријантност} значи да дескриптор исте физичке структуре остаје исти или довољно сличан после одређених промена снимања: померања, ротације, размере и умерене промене осветљења. \textbf{Препознатљивост}, односно \en{distinctiveness}, значи да различите структуре имају довољно различите дескрипторе да их можемо разликовати.

При упаривању две слике исте сцене камера ретко има исти положај, размеру и експозицију. Без робусности на те промене, исто обележје не би било пронађено као најближи сусед свог пара.
\[
\text{иста структура, друга слика}\ \Rightarrow\ d_1\approx d_2.
\]
Ипак, инваријантност је увек према \textbf{наведеној групи промена}, не према свему. Потпуно уклањање информација може уништити препознатљивост. SIFT канонизује скалу и оријентацију и нормализује градијенте, али није савршено инваријантан на произвољну перспективу, заклањање или промену осветљења.
\izvor{Прочитати §7.1.2, \en{Feature descriptors}, и §7.1.3, \en{Feature matching}, из [S].}
\end{odgovor}

\begin{odgovor}{Које сопствене вредности матрице момента другог реда код Harris детектора указују на добро обележје за праћење? Графички илустровати.}
За локално суседство и тежине $w$:
\[
M=\sum_{x,y}w(x,y)
\begin{bmatrix}I_x^2&I_xI_y\\I_xI_y&I_y^2\end{bmatrix},\qquad
E(\Delta)\approx\Delta^\top M\Delta.
\]
Сопствене вредности $\lambda_1,\lambda_2\geq0$ мере промену изгледа при малом померању у два главна правца.
\begin{center}\begin{tikzpicture}[scale=.75]
\draw[arr](0,0)--(6,0) node[right]{$\lambda_1$};
\draw[arr](0,0)--(0,2.7) node[above]{$\lambda_2$};
\node[box] at (1,.35){равна област};\node[box] at (4.5,.35){ивица};
\node[box] at (1,2){ивица};\node[box,fill=green!15] at (4.5,2){угао: добар за праћење};
\end{tikzpicture}\end{center}
Обе мале вредности означавају равну област; једна велика и друга мала ивицу; \textbf{обе велике} угао или богату текстуру. Тада су оба смера померања добро одређена. Harris користи $R=\det M-k(\operatorname{tr}M)^2$; критеријум \en{Shi--Tomasi} непосредно прати $\min(\lambda_1,\lambda_2)$. Велика само једна вредност није довољна.
\izvor{Прочитати §7.1.1, \en{Feature detectors}, и §7.1.5, \en{Feature tracking}, из [S].}
\end{odgovor}

\begin{odgovor}{Како детектор обележја, нпр. SIFT, учинити инваријантним на промене осветљења?}
За идеализовану промену $I'=aI+b$, где је $a>0$ константан контраст, а $b$ константан померај осветљења, просторни градијенти уклањају $b$, док нормализација њихових величина уклања фактор $a$:
\[
\nabla I'=a\nabla I,\qquad
\frac{\nabla I'}{\|\nabla I'\|}=\frac{\nabla I}{\|\nabla I\|}.
\]
Зато се у SIFT дескриптору користе оријентације градијента и нормализовани хистограми, уз ограничење великих компоненти и поновну нормализацију. То даје робусност на умерене афине промене интензитета.

Треба разликовати \textbf{детектор} од \textbf{дескриптора}. DoG детекција уклања константан адитивни померај, али се њена амплитуда множи са $a$: фиксни праг контраста није строго инваријантан на скалирање осветљења. За робуснију детекцију могу се нормализовати локални контраст или користити релативни прагови. Сенке, засићење и просторне промене осветљења нарушавају једноставан модел; инваријантност није потпуна.
\izvor{Прочитати §7.1.1, \en{Feature detectors}, и §7.1.2, \en{Feature descriptors}, из [S], делове о DoG/SIFT и нормализацији дескриптора.}
\end{odgovor}

\begin{odgovor}{Да ли су Harris углови инваријантни на ротацију? Каква је ситуација код SIFT-а?}
При ротацији слике за матрицу $Q$, идеална матрица локалних градијената трансформише се као
\[
M'=QMQ^\top.
\]
Сопствене вредности, детерминанта и траг остају исти, па и Harris одзив. Зато је детекција углова ротационо коваријантна: угао се детектује на одговарајућој ротираној позицији, а јачина одзива је инваријантна. У пракси узорковање и интерполација дају одступања; изотропно суседство подржава ову особину.

\textbf{Harris детектор сам не даје ротационо инваријантан дескриптор}. Поређење сирових неротираних исечака и даље је осетљиво на ротацију.

SIFT додељује кључној тачки доминантну оријентацију из локалног хистограма градијената. Положаји у исечку и оријентације градијената потом се изражавају релативно према њој. Тако дескриптор добија ротациону инваријантност, а избор скале омогућава и робусност на промену размере. То не значи инваријантност на све геометријске промене.
\izvor{Прочитати §7.1.1 и §7.1.2 из [S], Harris детектор и додељивање оријентације у SIFT-у.}
\end{odgovor}

\begin{odgovor}{Како се формира вектор SIFT дескриптора?}
Око кључне тачке узмемо исечак на њеној скали и оријентишемо координате према додељеној доминантној оријентацији. Исечак поделимо на $4\times4$ просторне ћелије. У свакој акумулирамо хистограм са 8 оријентационих бина:
\[
4\cdot4\cdot8=128\quad\text{компоненти}.
\]
\begin{center}\begin{tikzpicture}[scale=.45]
\draw[step=1,blue,thick](0,0) grid(4,4);
\foreach \x in {.5,1.5,2.5,3.5}{\foreach \y in {.5,1.5,2.5,3.5}{\draw[arr](\x-.2,\y-.15)--(\x+.2,\y+.15);}}
\node[anchor=west] at (5,2){16 ћелија $\times$ 8 смерова $\to$ 128 вредности};
\end{tikzpicture}\end{center}
Гласови су пондерисани величином градијента и Гаусовом тежином према центру исечка. Интерполација између просторних ћелија и оријентационих бина смањује осетљивост на мале помаке. Хистограми се спајају у један вектор, L2 нормализују, велике компоненте ограничавају, типично на 0.2, и поново нормализују.

Просторне ћелије задржавају груб распоред обележја, док оријентациони хистограми и нормализација дају робусност на мање деформације и промене осветљења.
\izvor{Прочитати §7.1.2, \en{Feature descriptors}, из [S], део о SIFT дескриптору.}
\end{odgovor}

\begin{odgovor}{Како се RANSAC користи за спајање слика у панораму? Скицирати поступак.}
Прво детектујемо и опишемо обележја у преклапајућим сликама, па формирамо кандидате за упаривање. Неки парови су погрешни, зато се хомографија процењује робусно.
\lanac{Слике $\to$ обележја и парови $\to$ RANSAC $\to$ хомографија $H$ $\to$ \en{warping} $\to$ \en{blending}}
У једној RANSAC итерацији насумично бирамо 4 пара у општем положају, процењујемо $H$ и за све парове рачунамо грешку, нпр. растојање пројекције:
\[
e_i=\left\|x'_i-\pi(H\widetilde x_i)\right\|.
\]
Парови испод прага су \en{inliers}. Задржавамо модел са добрим консензусом и поново га процењујемо из свих његових исправних парова. Број итерација може се прилагодити односу исправних парова.

Хомографија је одговарајући модел за равну сцену или чисту ротацију камере око центра. За општу 3D сцену са транслацијом настаје паралакса, па једна хомографија није довољна. За више слика могу се додати глобално поравнање и корекција експозиције.
\izvor{Прочитати §8.1.4, \en{Robust least squares and RANSAC}, §8.2, \en{Image stitching}, и §8.4, \en{Compositing}, из [S].}
\end{odgovor}

\begin{odgovor}{Зашто се при PCA редукцији димензија прво одузима средња вредност тачака?}
За узорке $x_i\in\mathbb R^d$ рачунамо $\mu=\frac1N\sum_i x_i$ и центриране векторе $z_i=x_i-\mu$. PCA тражи правце највеће \textbf{варијансе око средње вредности}:
\[
C=\frac1N\sum_i z_iz_i^\top.
\]
Сопствени вектори $C$ са највећим сопственим вредностима задају главне компоненте. Координате у простору од $k$ компоненти су $a_i=U_k^\top(x_i-\mu)$.

Без центрирања, $\frac1N\sum_i x_ix_i^\top=C+\mu\mu^\top$ је други момент око исходишта, а не коваријанса. Велика удаљеност центра података од нуле могла би да доминира првом компонентом, иако не описује њихову стварну варијабилност. Центрирање помера облак тачака у исходиште, па правци зависе од облика облака, а не од његовог апсолутног положаја.

Одузимање средње вредности није исто што и стандардизација сваке компоненте по варијанси; стандардизација је додатни избор који мења PCA геометрију.
\izvor{Прочитати §5.2.3, \en{Principal component analysis}, из [S].}
\end{odgovor}

\begin{odgovor}{Како се PCA може искористити за детекцију лица на слици?}
Од поравнатих тренинг слика лица исте величине формирамо векторе пиксела, средње лице $\mu$ и PCA базу $U_k$. Њени вектори, приказани као слике, називају се \en{eigenfaces}.

Кандидатски прозор $x$ на новој слици, после уједначавања величине и осветљења, пројектујемо у простор лица:
\[
a=U_k^\top(x-\mu),\qquad \widehat x=\mu+U_ka,\qquad
e=\|x-\widehat x\|^2.
\]
Мала реконструкциона грешка указује да је прозор близу наученог простора лица. Слика се претражује на више позиција и скала, а прихваћени прозори спајају или филтрирају. Може се додати модел расподеле PCA коефицијената и обучити класификатор на лицима и не-лицима.

\textbf{Сам PCA није довољан детектор}: не-лице такође може имати малу грешку, а лице под другом позом или осветљењем велику. За препознавање идентитета пореде се коефицијенти већ издвојених лица; то је други задатак од детекције.
\izvor{Прочитати §5.2.3, \en{Principal component analysis}, §6.2.4, \en{Face recognition}, и §6.3.1, \en{Face detection}, из [S].}
\end{odgovor}

\begin{odgovor}{Која је разлика између генеративних и дискриминативних метода класификације?}
Генеративни приступ моделује како подаци настају у свакој класи: учи $p(x\mid y)$ и $p(y)$, односно заједничку расподелу $p(x,y)$. Класификација се добија Бајесовим правилом:
\[
p(y\mid x)=\frac{p(x\mid y)p(y)}{\sum_{y'}p(x\mid y')p(y')}.
\]
Примери су \en{Gaussian Bayes} и \en{naive Bayes}. Модел може описати или узорковати податке, али лоше претпоставке о њиховој расподели могу ослабити класификацију.

Дискриминативни приступ непосредно учи $p(y\mid x)$, као логистичка регресија, или границу и резултат одлучивања, као SVM. Не мора да објасни целу расподелу улазних података, већ само оно што је потребно за разликовање класа.

То није подела на „обучавано“ и „необучавано“: обе врсте могу користити означене примере. Генеративна класификација такође не значи обавезно генерисање слика GAN мрежом; овде је реч о начину моделирања расподела за класификацију.
\izvor{Прочитати §5.1.2, \en{Bayesian classification}, §5.1.3, \en{Logistic regression}, и §5.1.4, \en{Support vector machines}, из [S].}
\end{odgovor}

\begin{odgovor}{Како се објекти детектују методом клизајућег прозора?}
\en{Sliding window} испитује многе правоугаоне исечке слике. Прозор одређене величине помера се по слици са изабраним кораком. За сваки исечак издвојимо обележја и применимо класификатор „објекат / позадина“.
\begin{center}\begin{tikzpicture}[scale=.65]
\draw[fill=gray!10](0,0) rectangle(6,2.6);
\foreach \x in {.3,1.3,2.3,3.3}{\draw[blue,thick](\x,.5) rectangle(\x+1.5,2);}
\draw[arr](.5,.2)--(5,.2);
\node[box] at (10,1.3){Обележја $\to$ класификатор};
\draw[arr](6.2,1.3)--(7.6,1.3);
\end{tikzpicture}\end{center}
За различите величине објекта користи се пирамида слика или више размера прозора. Прозори чији резултат пређе праг постају кандидати. \en{Non-maximum suppression} задржава јаке детекције и уклања преклапајуће дупликате.

Мањи корак и више скала побољшавају покривеност, али повећавају цену. Број испитивања расте са бројем позиција и размера. Метод је једноставан, али без дељења израчунавања може бити спор и ограничен фиксним обликом прозора.
\izvor{Прочитати §6.3.1, \en{Face detection}, и §6.3.2, \en{Pedestrian detection}, из [S].}
\end{odgovor}

\begin{odgovor}{Како каскадни класификатор детектује лица и која је главна предност?}
Каскада примењује низ класификатора на сваки кандидатски прозор:
\lanac{Прозор $\to$ јефтина фаза 1 $\to$ фаза 2 $\to\cdots\to$ сложеније фазе $\to$ лице}
\lanac{\textit{Одбијање у било којој фази} $\Longrightarrow$ одмах прекид обраде прозора}
У Viola--Jones приступу фазе су засноване на \en{Haar-like} обележјима, брзо рачунатим помоћу интегралне слике, и \en{boosting}-у. Ране фазе користе мали број обележја и уклањају очигледну позадину; касније детаљније проверавају преостале кандидате.

\textbf{Главна предност је мала просечна цена}: већина прозора је позадина и одбацује се пре скупих фаза. Ране фазе се подешавају за висок \en{recall}, јер једном одбачено лице не може бити враћено. Ако је одзив фазе $r_j$, укупни одзив је производ условних одзива $\prod_jr_j$; чак и мали губици се акумулирају. Проценат лажних прихватања такође опада кроз каскаду.
\izvor{Прочитати §6.3.1, \en{Face detection}, из [S], Viola--Jones каскаду и \en{boosting}.}
\end{odgovor}

\begin{odgovor}{Шта су интегралне слике и како се користе при детекцији лица?}
Интегрална слика, односно \en{summed-area table}, чува збир пиксела у правоугаонику од исходишта до дате позиције:
\[
S(x,y)=\sum_{i\leq x,\ j\leq y}I(i,j).
\]
Може се израчунати у $O(HW)$ користећи
\[
S(x,y)=I(x,y)+S(x-1,y)+S(x,y-1)-S(x-1,y-1).
\]
Збир произвољног осно поравнатог правоугаоника $[x_1,x_2]\times[y_1,y_2]$ је
\[
S(x_2,y_2)-S(x_1-1,y_2)-S(x_2,y_1-1)+S(x_1-1,y_1-1),
\]
уз нулу за позиције изван горње и леве границе.

Зато су потребна само четири приступа табели, независно од величине правоугаоника. \en{Haar-like} обележје је разлика збира у светлим и тамним правоугаоницима, па се много таквих обележја брзо рачуна у сваком прозору Viola--Jones детектора. Интегрална слика сама не препознаје лице: она само убрзава израчунавање обележја.
\izvor{Прочитати §3.2, \en{Linear filtering}, део о \en{summed-area tables}, и §6.3.1, \en{Face detection}, из [S].}
\end{odgovor}

\begin{odgovor}{Како се SVM користи када подаци нису линеарно сепарабилни?}
Постоје два различита решења. \textbf{Soft-margin SVM} дозвољава нарушавање маргине и неке погрешне класификације:
\[
\min_{w,b,\xi}\frac12\|w\|^2+C\sum_i\xi_i,\qquad
y_i(w^\top x_i+b)\geq1-\xi_i,\quad\xi_i\geq0.
\]
$C$ контролише компромис између широке маргине и кажњавања грешака. Овај приступ је користан за шум и преклапајуће класе, али граница остаје линеарна.

За нелинеарну границу користи се пресликавање $\phi(x)$ и \en{kernel trick}:
\[
k(x_i,x_j)=\phi(x_i)^\top\phi(x_j),\qquad
f(x)=\sum_i\alpha_i y_i k(x_i,x)+b.
\]
На пример, RBF језгро $k(x,x')=\exp(-\gamma\|x-x'\|^2)$ омогућава сложену границу без експлицитног рачунања високодимензионалног $\phi$. Ова два механизма могу се комбиновати. Параметри $C$, $\gamma$ и избор језгра бирају се на валидационим подацима; превише флексибилна граница може се преобучити.
\izvor{Прочитати §5.1.4, \en{Support vector machines}, из [S], \en{soft margin} и језгра.}
\end{odgovor}

\begin{odgovor}{Како HOG и линеарни SVM детектују особе на слици?}
\en{Histogram of oriented gradients} — HOG — описује локалну структуру обриса. Израчунамо градијенте, поделимо прозор на ћелије и у свакој формирамо хистограм оријентација, пондерисан величином градијента. Често се користе неусмерене оријентације $0^\circ$--$180^\circ$.

Суседне ћелије групишу се у преклапајуће блокове; њихови хистограми се нормализују, нпр. поступком \en{L2-Hys}. Нормализација ублажава утицај осветљења и контраста. Сви блокови се спајају у вектор $h$.
\lanac{Прозор $\to$ градијенти $\to$ ћелије и хистограми $\to$ нормализовани блокови $\to$ SVM}
Линеарни SVM, обучен на прозорима са особама и позадином, рачуна $s=w^\top h+b$. Применимо га клизањем прозора на више скала, прагом издвојимо кандидате и уклонимо дупликате помоћу \en{NMS}. \en{Hard-negative mining} допуњава тренинг лажним позитивима.

HOG задржава груб просторни распоред градијената, што описује облик особе, без ослањања само на њену боју.
\izvor{Прочитати §6.3.2, \en{Pedestrian detection}, и §5.1.4, \en{Support vector machines}, из [S].}
\end{odgovor}

\begin{odgovor}{Како се bag of visual words користи за претрагу слика?}
На тренинг сликама издвојимо локалне дескрипторе, нпр. SIFT, и групишемо их методом \en{k-means}. Центри кластера чине визуелни речник. Сваки дескриптор нове слике додељујемо најближој визуелној речи, па слику представљамо хистограмом њихових појављивања.
\lanac{Дескриптори $\to$ визуелни речник $\to$ хистограм речи $\to$ рангирање сличних слика}
За упитну слику рачунамо исту репрезентацију и поредимо је са базом, нпр. косинусном сличношћу. \en{TF--IDF} пондерисање смањује значај речи које се јављају у скоро свим сликама, а повећава значај специфичнијих речи.

Предност је компактна репрезентација погодна за индексирање великих база. Основни \en{bag-of-words} губи просторни распоред: исти хистограм могу имати различите сцене. Зато се најбољи кандидати могу додатно проверити геометријским упаривањем, нпр. RANSAC-ом. Визуелна реч није нужно семантички објекат, већ кластер локалних дескриптора.
\izvor{Прочитати §7.1.4, \en{Large-scale matching and retrieval}, §6.2.1, \en{Feature-based methods}, и §6.2.3 из [S].}
\end{odgovor}

\begin{odgovor}{Објаснити употребу инверзног индекса за визуелне речи.}
Инверзни индекс за сваку визуелну реч чува листу слика у којима се јавља, уз број појављивања или тежину:
\begin{center}\begin{tabular}{ll}\toprule
Визуелна реч & Листа појављивања — \en{posting list}\\\midrule
$w_1$ & $(I_2,0.4), (I_7,0.8)$\\
$w_2$ & $(I_1,0.6), (I_7,0.2), (I_9,0.3)$\\
$w_3$ & $(I_2,0.7), (I_9,0.5)$\\\bottomrule
\end{tabular}\end{center}
Када упит садржи $w_1$ и $w_3$, пролазимо само кроз њихове листе и акумулирамо резултат за слике $I_2,I_7,I_9$. Не морамо поредити упит са сваким дескриптором сваке слике.

Често се сумирају производи \en{TF--IDF} тежина, па добијени резултати нормализују и рангирају. Цена зависи од укупне дужине листа речи упита, а не нужно од величине целе базе. Веома честе речи дају дуге листе и слабу дискриминацију, па се умањују или занемарују. Геометријска провера може се применити на најбоље рангиране слике.
\izvor{Прочитати §7.1.4, \en{Large-scale matching and retrieval}, и §6.2.3, \en{Application: Visual similarity search}, из [S].}
\end{odgovor}

\begin{odgovor}{Описати модел вештачког неурона и ReLU. Зашто је ова активација заступљена у дубоким мрежама?}
Неурон рачуна пондерисани збир улаза са померајем, па примењује нелинеарну активацију:
\[
a=w^\top x+b,\qquad h=\varphi(a),\qquad \operatorname{ReLU}(a)=\max(0,a).
\]
\begin{center}\begin{tikzpicture}[scale=.6]
\draw[arr](-3,0)--(3.5,0) node[right]{$a$};\draw[arr](0,-.4)--(0,3.2) node[above]{$h$};
\draw[blue,very thick](-3,0)--(0,0)--(3,3);
\end{tikzpicture}\end{center}
\en{ReLU} је рачунски једноставна и за позитивне улазе има извод 1, па не засићује као сигмоид. То олакшава проток градијента у дубоким мрежама. Негативни улази дају нулу, што подстиче ретке активације; нелинеарност омогућава сложене границе одлучивања.

За негативан улаз извод је 0, па неурон може постати \en{dead ReLU} и престати да учи. У нули извод није јединствен и имплементација бира конвенцију, обично 0. \en{Leaky ReLU} уводи мали негативни нагиб. ReLU сама не гарантује одсуство нестајућих градијената.
\izvor{Прочитати §5.3.1, \en{Weights and layers}, и §5.3.2, \en{Activation functions}, из [S].}
\end{odgovor}

\begin{odgovor}{Шта је overfitting при обучавању неуронских мрежа и како се избегава?}
\en{Overfitting} значи да модел добро описује тренинг примере, укључујући њихов шум и случајне особине, али слабо генерализује на нове податке. Типичан знак је опадање тренинг грешке уз раст валидационе грешке.

Поступци за ублажавање су L2 регуларизација или одговарајући \en{weight decay}, \en{dropout}, \en{data augmentation}, мањи капацитет модела, више разноврсних података и \en{transfer learning}. \en{Early stopping} задржава тежине из корака са најбољим валидационим резултатом, уместо последњег корака.

Параметри се уче на тренинг скупу, а хиперпараметри и тренутак заустављања бирају на валидационом. Тест скуп служи коначној процени, не подешавању. Потребно је спречити \en{data leakage}, нпр. да готово исте слике или исти субјекти заврше у тренинг и тест скупу. Регуларизација није бесплатна: ако је превише јака или модел премали, настаје \en{underfitting}.
\izvor{Прочитати §4.1.2, \en{Overfitting and underfitting}, §5.3.3, \en{Regularization and normalization}, и §5.3.6 из [S].}
\end{odgovor}

\begin{odgovor}{Зашто су дубоке CNN мреже ефикасне у различитим задацима рачунарског вида?}
CNN користи структуру слике: локални филтри обрађују суседне пикселе, а исте тежине деле се преко просторних позиција. То смањује број параметара у односу на потпуно повезану обраду и омогућава да се исти образац препозна на различитим местима.

Дубина гради хијерархију обележја: од ивица и текстура до делова објеката и семантичких образаца. Обележја и завршна одлука уче се заједно, \en{end-to-end}, уместо да се сва обележја ручно пројектују. Велики скупови података, GPU израчунавање и \en{pretraining} додатно подржавају њихову успешност.

Конволуција је, под идеалним условима, транслационо \textbf{еквиваријантна}: померање улаза помера и мапу обележја. Потпуно инваријантан излаз не следи аутоматски; \en{pooling}, агрегација и обучавање могу дати приближну инваријантност. CNN мреже и даље могу грешити при промени домена, непознатим условима и намерним пертурбацијама.
\izvor{Прочитати §5.4, \en{Convolutional neural networks}, и §6.2.2, \en{Deep networks}, из [S].}
\end{odgovor}

\begin{odgovor}{Описати структуру типичне конволуционе неуронске мреже.}
\lanac{Слика $\to$ [конволуција $\to$ нормализација $\to$ активација] $\to$ смањивање резолуције}
\lanac{$\to$ дубљи блокови $\to$ \en{global average pooling} / \en{flatten} $\to$ завршни слој}
Конволуциони слојеви издвајају локална обележја, а активације, нпр. \en{ReLU}, уводе нелинеарност. Нормализација често олакшава обучавање. \en{Pooling} или конволуције са кораком већим од 1 смањују висину и ширину мапа; број канала се често повећава у дубљим блоковима.

Код класификације се просторне мапе агрегирају, а последњи слој даје резултате за класе. \en{Softmax} може их претворити у расподелу вероватноћа. Код детекције или сегментације завршна \en{head} даје оквире или просторне маске уместо једне ознаке.

Ово је општи образац, не обавезан редослед сваке архитектуре. Савремене мреже могу имати резидуалне везе и више грана; \en{global average pooling} често замењује велике потпуно повезане слојеве.
\izvor{Прочитати §5.4, §5.4.1, \en{Pooling and unpooling}, и §5.4.3, \en{Network architectures}, из [S].}
\end{odgovor}

\begin{odgovor}{Која је улога и како се формира излаз конволуционог слоја? Илустровати примером.}
Филтар се помера преко слике. На свакој позицији рачуна скаларни производ тежина са локалним прозором преко свих улазних канала и додаје \en{bias}. Један филтар даје једну излазну мапу.
\[
y_{i,j,k}=b_k+\sum_{u,v,c}w_{u,v,c,k}\,x_{iS+u-P,jS+v-P,c}.
\]
За прозор и језгро
\[
\begin{bmatrix}1&2\\3&4\end{bmatrix},\qquad
\begin{bmatrix}1&0\\0&-1\end{bmatrix}
\]
скаларни производ је $1-4=-3$, ако је $b=0$. Исти филтар на другом положају даје други елемент мапе. Локална повезаност и дељење тежина издвајају обрасце уз мало параметара.

За улаз $H\times W\times C$, језгро $K\times K$, корак $S$, попуну $P$ и без дилатације:
\[
H_{\mathrm{out}}=\left\lfloor\frac{H+2P-K}{S}\right\rfloor+1.
\]
Број излазних канала једнак је броју филтара. У библиотекама операција обично представља \en{cross-correlation}, без обртања језгра.
\izvor{Прочитати §5.4, \en{Convolutional neural networks}, и §3.2, \en{Linear filtering}, из [S].}
\end{odgovor}

\begin{odgovor}{Која је улога и како се формира излаз pooling слоја? Илустровати примером.}
\en{Pooling} сажима вредности у локалном прозору сваког канала независно. \en{Max pooling} бира максимум, а \en{average pooling} средњу вредност. Типично нема обучиве параметре.
\[
\begin{bmatrix}
1&3&2&4\\
0&2&5&1\\
7&6&0&2\\
3&1&4&8
\end{bmatrix}
\xrightarrow[\text{корак }2]{\textit{max pooling }2\times2}
\begin{bmatrix}3&5\\7&8\end{bmatrix}.
\]
Мања просторна резолуција смањује цену наредних слојева и може ублажити утицај малих померања, али губи прецизан положај и ситне детаље. Број канала остаје непромењен. Димензије излаза рачунају се као за локални филтар, уз задате кораке, попуну и конвенцију заокруживања.

У \en{backpropagation}-у максимум прослеђује градијент позицији која је дала највећу вредност, док просек градијент дели по целом прозору. Pooling не даје потпуну инваријантност на произвољна померања.
\izvor{Прочитати §5.4.1, \en{Pooling and unpooling}, из [S].}
\end{odgovor}

\begin{odgovor}{Где се 3D тачка $(X,Y,Z)$ у простору камере пресликава при перспективној пројекцији? Скицирати и написати $u,v$.}
За виртуелну раван на растојању $f$ испред центра камере, сличност троуглова даје физичке координате $x=fX/Z$, $y=fY/Z$.
\projsketch
У пикселским координатама, без дисторзије:
\[
u=f_x\frac XZ+s\frac YZ+c_x,\qquad
v=f_y\frac YZ+c_y.
\]
$f_x,f_y$ су жижне даљине изражене у пикселима, $(c_x,c_y)$ главна тачка, а $s$ параметар искошења оса, обично 0. Скица приказује пресек у равни $XZ$; за $Y$ важи исти однос.

Дељење са $Z$ значи да даље тачке, при истом $X,Y$, имају мањи отклон од главне тачке. Све тачке на истом зраку из центра имају исту пројекцију, па једна слика сама не одређује њихову дубину. За модел са физичком равни иза центра јавља се знак минус.
\izvor{Прочитати §2.1.4, \en{3D to 2D projections}, из [S].}
\end{odgovor}

\begin{odgovor}{Илустровати pinhole модел камере. Који параметри одређују intrinsic и extrinsic матрицу?}
\en{Pinhole camera} пројектује тачку на пресек зрака кроз центар камере и раван слике:
\projsketch
\[
\widetilde x\sim P\widetilde X_w,\qquad P=K[R\mid t],\qquad
K=\begin{bmatrix}f_x&s&c_x\\0&f_y&c_y\\0&0&1\end{bmatrix}.
\]
\textbf{Intrinsic} параметри $f_x,f_y,s,c_x,c_y$ описују пројекцију из система камере у пикселе: жижну даљину и размере пиксела, искошење оса и главну тачку.

\textbf{Extrinsic} параметри описују положај и оријентацију камере у односу на свет. Овде $X_c=RX_w+t$: $R$ је ротација, а $t$ транслација у пресликавању свет $\to$ камера. Центар камере у свету је $C=-R^\top t$; $t$ није директно њен положај. Дисторзија сочива захтева додатне параметре и није садржана у идеалној матрици $K$.
\izvor{Прочитати §2.1.4, \en{3D to 2D projections}, и §2.1.5, \en{Lens distortions}, из [S].}
\end{odgovor}

\begin{odgovor}{Зашто паралелне линије конвергирају у vanishing point на слици?}
Паралелне 3D праве деле смер $d$. Тачке на њима можемо записати као $X(\lambda)=X_0+\lambda d$. Како $\lambda\to\infty$, у перспективној пројекцији доминира смер, а утицај почетне тачке нестаје. За смер изражен у светском систему:
\[
v_{\mathrm{hom}}\sim KRd.
\]
Зато пројекције паралелних права пролазе кроз исту тачку недогледа, односно \en{vanishing point}. Она зависи од смера и оријентације камере, а не од удаљености праве или транслације камере.
\begin{center}\begin{tikzpicture}[scale=.7]
\draw[thick](-4,-1.2)--(0,1.4); \draw[thick](4,-1.2)--(0,1.4);
\draw[thick](-2,-1.2)--(0,1.4); \draw[thick](2,-1.2)--(0,1.4);
\fill[red](0,1.4) circle(3pt) node[above]{$v$};
\node at (0,-1.6){пројекције паралелних 3D права};
\end{tikzpicture}\end{center}
Ако је смер паралелан равни слике, хомогена координата дубине је нула: тачка недогледа је у бесконачности, па пројектоване праве остају паралелне. Различити 3D смерови уопштено имају различите тачке недогледа.
\izvor{Прочитати §2.1.4, §7.4.3, \en{Vanishing points}, и §11.1.1 из [S].}
\end{odgovor}

\begin{odgovor}{Колико је најмање 3D--2D парова потребно за параметре камере? Како се процењују са више парова?}
За \textbf{општу непознату пројекциону матрицу} $P\in\mathbb R^{3\times4}$ постоји 11 степени слободе: 12 елемената, до заједничке размере. Један 3D--2D пар даје две независне једначине, па је потребно најмање \textbf{6 парова у недегенерисаном просторном распореду}. Сам број није довољан: све тачке у једној равни не одређују општу $P$.
\[
\widetilde x_i\times P\widetilde X_i=0
\quad\Longrightarrow\quad Ap=0.
\]
\en{Direct Linear Transform} — DLT — нормализује координате и гради матрицу $A$. Са више парова, десни сингуларни вектор за најмању сингуларну вредност даје линеарну процену. Потом се нелинеарно минимизује \textbf{репројекциона грешка}; RANSAC може одбацити погрешне парове.

Ако је $K$ већ позната и тражи се само положај и оријентација, то је PnP проблем: P3P користи 3 пара, али може дати више решења, па се обично додатним паром разрешава двосмисленост. Планарна калибрација из више погледа је засебан случај.
\izvor{Прочитати §11.2.1, \en{Linear algorithms}, §11.2.2, \en{Iterative non-linear algorithms}, и §11.1, \en{Geometric intrinsic calibration}, из [S].}
\end{odgovor}

\begin{odgovor}{Шта је афина трансформација и колико парова тачака је потребно? Графички илустровати.}
За 2D тачке афина трансформација је
\[
\begin{bmatrix}u\\v\end{bmatrix}
=\begin{bmatrix}a_{11}&a_{12}\\a_{21}&a_{22}\end{bmatrix}
\begin{bmatrix}x\\y\end{bmatrix}
+\begin{bmatrix}t_x\\t_y\end{bmatrix}.
\]
Има 6 параметара: линеарни део може комбиновати ротацију, скалирање, смицање и рефлексију, а $t$ је транслација. Чува праве, паралелност и односе растојања на истој правој, али не мора чувати углове и дужине.
\begin{center}\begin{tikzpicture}[scale=.7]
\draw[blue,thick](0,0) rectangle(2,1.5);
\draw[arr](2.5,.75)--(4,.75);
\draw[red,thick](5,0)--(8,0)--(9,1.5)--(6,1.5)--cycle;
\node at (1,-.4){правоугаоник};\node at (7,-.4){паралелограм};
\end{tikzpicture}\end{center}
Један пар даје две једначине, па су минимално потребна \textbf{3 пара са неколинеарним изворним тачкама}. За недегенерисано, инвертибилно пресликавање и циљни троугао мора имати ненулту површину. Са више парова користи се метод најмањих квадрата, а уз погрешне парове робусна процена.
\izvor{Прочитати §2.1.1, \en{2D transformations}, и §8.1.1, \en{2D alignment using least squares}, из [S].}
\end{odgovor}

\begin{odgovor}{Шта је хомографија и колико парова тачака је потребно? Графички илустровати.}
Хомографија је инвертибилна пројективна трансформација између две равни:
\[
\widetilde x'\sim H\widetilde x,\qquad H\in\mathbb R^{3\times3},
\]
\[
u=\frac{h_{11}x+h_{12}y+h_{13}}{h_{31}x+h_{32}y+h_{33}},\qquad
v=\frac{h_{21}x+h_{22}y+h_{23}}{h_{31}x+h_{32}y+h_{33}}.
\]
Има 8 степени слободе, јер заједничка размера матрице није битна. Чува праве и инциденцију, али уопштено не паралелност, углове или дужине.
\begin{center}\begin{tikzpicture}[scale=.65]
\draw[blue,thick](0,0) rectangle(2,1.5);\draw[arr](2.5,.75)--(4,.75);
\draw[red,thick](5,0)--(8,0)--(7.4,1.5)--(5.5,1.5)--cycle;
\node at (1,-.4){раван};\node at (6.5,-.4){перспективни поглед};
\end{tikzpicture}\end{center}
Потребна су \textbf{4 пара у општем положају}: међу четири тачке ни у једној слици не треба да постоје три колинеарне. Са више парова користи се нормализовани DLT и репројекционо прецизирање. Хомографија описује погледе равне сцене или чисту ротацију камере; није општи модел две слике произвољне 3D сцене са паралаксом.
\izvor{Прочитати §2.1.1, \en{2D transformations}, и §8.1.1, \en{2D alignment using least squares}, из [S].}
\end{odgovor}

\begin{odgovor}{За две камере и просторну тачку $X$, приказати епиполарне линије на равнима слика.}
Центри $C_1,C_2$ и тачка $X$ одређују епиполарну раван $\Pi$. Њени пресеци са равнима две слике су епиполарне праве $\ell$ и $\ell'$. Пројекције $x$ и $x'$ леже на тим правама.
\episketch
На скици је простор приказан шематски: сиво су равни слика, плавкасти троугао је део епиполарне равни, а обојене линије њени пресеци. Свака епиполарна линија пролази кроз одговарајући епипол, пројекцију другог центра камере.

За познат пиксел $x$ не знамо дубину, али одговарајућа 3D тачка мора лежати на његовом зраку. Њена пројекција у другој камери зато мора бити на $\ell'$:
\[
\ell'=F\widetilde x,\qquad \ell=F^\top\widetilde x',\qquad
\widetilde x'^\top F\widetilde x=0.
\]
Епиполарно ограничење своди просторну претрагу пара са целе слике на једну линију.
\izvor{Прочитати §12.1, \en{Epipolar geometry}, и §11.3, \en{Two-frame structure from motion}, из [S].}
\end{odgovor}

\begin{odgovor}{Шта је епипол? Где је код кретања дуж осе камере, а где код кретања управно на њу? Илустровати.}
Епипол је пројекција центра једне камере на слику друге, односно пресек основице са њеном пројективном равни. Кроз њега пролазе све епиполарне линије. За фундаменталну матрицу важи $Fe=0$ и $F^\top e'=0$.

За \textbf{чисту транслацију без промене оријентације}, дуж оптичке осе, епиполи су у главним тачкама слика. За чисту транслацију паралелну равни слике, односно управну на оптичку осу, епиполи су у бесконачности: епиполарне линије су паралелне.
\begin{center}\begin{tikzpicture}[scale=.7]
\draw(0,0) rectangle(4,2);
\foreach \x/\y in {0/0,0/2,4/0,4/2,0/1,4/1}{\draw[blue](2,1)--(\x,\y);}
\fill[red](2,1) circle(3pt) node[above]{$e$};
\draw(6,0) rectangle(10,2);
\foreach \y in {.4,1,1.6}{\draw[arr,blue](6,\y)--(10.5,\y);}
\node at (2,-.5){осно кретање};\node at (8,-.5){бочно: $e$ у бесконачности};
\end{tikzpicture}\end{center}
Ако постоји и ротација, положај епипола се мора одредити у сваком систему камере; наведене једноставне скице не важе аутоматски.
\izvor{Прочитати §12.1, \en{Epipolar geometry}, и §11.3.2, \en{Special motions and structures}, из [S].}
\end{odgovor}

\begin{odgovor}{Која је разлика између есенцијалне и фундаменталне матрице?}
Есенцијална матрица $E$ описује епиполарну геометрију у \textbf{нормализованим координатама калибрисаних камера}. За $X_2=RX_1+t$:
\[
E=[t]_\times R,\qquad \widehat x'^\top E\widehat x=0,
\]
где $[t]_\times a=t\times a$.

Фундаментална матрица $F$ делује над \textbf{пикселским координатама} и обухвата и унутрашње параметре:
\[
\widetilde x'^\top F\widetilde x=0,\qquad
F=K_2^{-\top}EK_1^{-1},\qquad E=K_2^\top FK_1.
\]
$F$ има ранг 2 и 7 степени слободе. $E$ има ранг 2, две једнаке ненулте сингуларне вредности и 5 степени слободе: три за ротацију и два за смер транслације. Апсолутна размера транслације не може се добити само из овог односа две слике.

Обе матрице дају ограничење за упаривање; не представљају хомографију која сваком пикселу непосредно додељује јединствен пиксел у другој слици.
\izvor{Прочитати §11.3.1, \en{Eight, seven, and five-point algorithms}, и §12.1, \en{Epipolar geometry}, из [S].}
\end{odgovor}

\begin{odgovor}{Какав је однос диспаритета и даљине тачке у калибрисаном стерео систему? Графички илустровати.}
За ректификоване камере са паралелним осама, основном раздаљином $B$ и истом жижном даљином $f_x$ у пикселима, користимо координате релативне главним тачкама:
\[
\bar u_L=f_xX/Z,\qquad \bar u_R=f_x(X-B)/Z,
\]
\[
d=\bar u_L-\bar u_R=\frac{f_xB}{Z},\qquad \boxed{Z=\frac{f_xB}{d}}.
\]
\begin{center}\begin{tikzpicture}[scale=.65]
\fill(0,0) circle(2pt) node[below]{$C_L$};\fill(6,0) circle(2pt) node[below]{$C_R$};
\fill(2.5,3.5) circle(2pt) node[above]{$X,Z$};
\draw[thick](0,0)--(2.5,3.5)--(6,0);
\draw[blue](-.5,1)--(1.8,1);\draw[blue](4.3,1)--(6.5,1);
\draw[dashed](0,0)--(0,1.7);\draw[dashed](6,0)--(6,1.7);
\fill[red](.714,1) circle(3pt);\fill[red](5,1) circle(3pt);
\draw[<->](0,-.6)--node[below]{$B$}(6,-.6);
\node at (3,1){равни $Z=f$};
\end{tikzpicture}\end{center}
Диспаритет је \textbf{обрнуто пропорционалан дубини}: ближи објекти имају већи помак, а врло далеки мали. Грешка дубине приближно расте као $|\delta Z|\approx Z^2|\delta d|/(f_xB)$. Ако главне тачке нису исте, сирово $u_L-u_R$ мора се кориговати њиховом разликом.
\izvor{Прочитати §12.1, \en{Epipolar geometry}, §12.1.1, \en{Rectification}, и увод поглавља 12 из [S].}
\end{odgovor}

\begin{odgovor}{Које ограничење се користи за налажење упарених пиксела у стерео сликама? Графички илустровати.}
Главно геометријско ограничење је \textbf{епиполарно ограничење}: пар пиксела мора задовољити $\widetilde x'^\top F\widetilde x=0$. За дати $x$ кандидат $x'$ тражимо само на $\ell'=F\widetilde x$.

После ректификације епиполарне линије постају хоризонталне, па парови имају исти ред, $v_L=v_R$:
\begin{center}\begin{tikzpicture}[scale=.7]
\draw(0,0) rectangle(4,2);\draw(6,0) rectangle(10,2);
\draw[blue,thick](0,1.2)--(4,1.2);\draw[blue,thick](6,1.2)--(10,1.2);
\fill[red](2.8,1.2) circle(3pt) node[above]{$x$};
\fill[red](7.5,1.2) circle(3pt) node[above]{$x'$};
\draw[arr](6.2,.7)--(9.8,.7);
\node at (2,-.4){лева слика};\node at (8,-.4){претрага истог реда};
\end{tikzpicture}\end{center}
Геометрија не бира јединствен пар; потребна је и сличност локалног изгледа, нпр. SSD, корелација или дескриптор. Додатни услови могу бити дозвољен распон диспаритета, глаткост унутар површине и лево-десна доследност. Заклоњени пиксели немају видљив пар, а понављајуће текстуре могу створити двосмисленост.
\izvor{Прочитати §12.1, §12.1.1, §12.3, \en{Dense correspondence}, и §12.4, \en{Local methods}, из [S].}
\end{odgovor}

\begin{odgovor}{Које претпоставке користи Lucas--Kanade оптички ток и како се примењује када оне не важе?}
Основне претпоставке су очување осветљености праћене тачке, мало померање које дозвољава линеаризацију и приближно исто кретање у локалном прозору:
\[
I(x+u,y+v,t+1)\approx I(x,y,t)
\quad\Longrightarrow\quad I_xu+I_yv+I_t=0.
\]
Један пиксел даје једну једначину за два непозната; \en{Lucas--Kanade} спаја једначине суседства и решава пондерисане најмање квадрате:
\[
\left(\sum w
\begin{bmatrix}I_x^2&I_xI_y\\I_xI_y&I_y^2\end{bmatrix}\right)
\begin{bmatrix}u\\v\end{bmatrix}
=-\sum wI_t\begin{bmatrix}I_x\\I_y\end{bmatrix}.
\]
За велика померања користи се пирамида \en{coarse-to-fine}, уз итеративно \en{warping} поравнање и поновну линеаризацију. За промене осветљења могу се процењивати локални контраст и померај интензитета или користити робуснија обележја. Робусна грешка ублажава утицај заклањања и погрешних пиксела.

Потребна је текстура у бар два смера. Пирамиде не решавају све проблеме: прозор преко границе два различита кретања и даље нарушава претпоставку локално истог тока.
\izvor{Прочитати §9.1, \en{Translational alignment}, §9.1.1, \en{Hierarchical motion estimation}, §9.1.3 и §9.3, \en{Optical flow}, из [S].}
\end{odgovor}

\begin{odgovor}{Шта је проблем апертуре при одређивању правца кретања? Илустровати примером.}
Кроз мало суседство које садржи само праву ивицу можемо одредити компоненту кретања нормалну на њу, али не и компоненту дуж ивице. Једначина
\[
I_xu+I_yv=-I_t
\]
ограничава скаларни производ брзине и градијента; не одређује цео вектор брзине.
\begin{center}\begin{tikzpicture}[scale=.75]
\draw[thick](-1,-1) rectangle(2,2);\draw[very thick,blue](0,-1)--(0,2);
\draw[arr,red](0,.5)--(1.5,.5) node[right]{нормална компонента};
\draw[arr](0,.5)--(1.5,1.7);\draw[arr](0,.5)--(1.5,-.7);
\node[font=\small] at (.4,-1.5){све три брзине имају исту хоризонталну компоненту};
\end{tikzpicture}\end{center}
За вертикалну ивицу $I_y\approx0$, па се може проценити $u$, док је $v$ неодређено. Различита стварна кретања зато имају исти локални визуелни ефекат.

Угао или текстурисана област садрже градијенте у више праваца и могу одредити обе компоненте. Проблем се ублажава ширим суседством, добрим избором обележја и додатним ограничењима глаткости или модела кретања. Равна област је још неодређенија: нема довољно информације ни о нормалном кретању.
\izvor{Прочитати §9.3, \en{Optical flow}, и §7.1.5, \en{Feature tracking}, из [S].}
\end{odgovor}

\begin{odgovor}{Која је веза Lucas--Kanade оптичког тока и Harris детектора углова?}
Обе методе користе исту локалну матрицу, односно \en{structure tensor}:
\[
M=\sum w
\begin{bmatrix}I_x^2&I_xI_y\\I_xI_y&I_y^2\end{bmatrix}.
\]
У Lucas--Kanade методу она је матрица нормалних једначина $M\delta=b$ за процену померања. Ако су обе сопствене вредности довољно велике, локална текстура одређује обе компоненте кретања и решење је стабилније. Једна мала вредност означава ивицу и проблем апертуре; обе мале равну област.

Harris користи детерминанту и траг исте матрице:
\[
R=\det M-k(\operatorname{tr}M)^2
\]
да пронађе углове. Зато су Harris углови природни кандидати за Lucas--Kanade праћење. \en{Shi--Tomasi} критеријум „добрих обележја за праћење“ непосредно тражи довољно велико $\min(\lambda_1,\lambda_2)$.

Иста матрица не значи исти задатак: Harris бира позиције у једној слици, а Lucas--Kanade користи и временску или међусликовну разлику да процени њихово померање.
\izvor{Прочитати §7.1.1, \en{Feature detectors}, §7.1.5, \en{Feature tracking}, и §9.1, \en{Translational alignment}, из [S].}
\end{odgovor}

\begin{odgovor}{Објаснити принцип дубинске камере Kinect-а.}
\textbf{Kinect v1} користи активно \en{structured light} мерење. Инфрацрвени пројектор емитује познат образац тачака, а IR камера снима његов изглед на сцени. У односу на референтни образац, положаји тачака се померају зависно од дубине.
\lanac{IR пројектор $\to$ познат образац на сцени $\to$ IR камера $\to$ упаривање и триангулација $\to$ дубина}
Пројектор и камера имају познат међусобни положај. Из помака и калибрације процењује се дубина, по принципу активног стереа; ближе површине дају другачији диспаритет од даљих. RGB камера обезбеђује боју, али није извор овог дубинског мерења.

\textbf{Kinect v2} користи \en{time-of-flight}: емитује модулисано IR светло и из фазног помераја повратног сигнала процењује време пута, односно удаљеност. Основни однос за повратни пут је $Z=c\Delta t/2$.

Зато „Kinect“ није један универзалан физички принцип. Јако спољно IR осветљење, заклањање, неповољна рефлексија и, код ToF-а, вишеструки путеви светла могу погоршати мерење.
\izvor{Прочитати §13.2, \en{3D scanning}, из [S], делове о \en{structured light}, Kinect-у и \en{time-of-flight} сензорима; за геометријску основу §12.1.}
\end{odgovor}

\begin{odgovor}{Објаснити Mean Shift и његову употребу за делове тела код Kinect-а.}
\en{Mean Shift} тражи модове, односно локалне максимуме процењене густине. Од текуће позиције $z_t$ рачуна се пондерисана средња вредност оближњих узорака и позиција помера ка њој:
\[
z_{t+1}=\frac{\sum_i w_i\,g(\|z_t-z_i\|^2/h^2)\,z_i}
{\sum_i w_i\,g(\|z_t-z_i\|^2/h^2)}.
\]
$h$ је \en{bandwidth}; $g$ је тежински профил повезан са језгром густине. Понављање води ка моду. При сегментацији слике узорци могу бити вектори положаја и боје; области које конвергирају ка истом моду групишу се.

У класичном Kinect поступку \textbf{класификацију пиксела у делове тела прво ради randomized decision forest}. Дубински пиксели се затим репројектују у 3D, пондеришу вероватноћама припадности делу и групишу \en{Mean Shift}-ом. Модови дају кандидате за положаје зглобова.
\lanac{Дубинска слика $\to$ класе делова тела $\to$ пондерисане 3D тачке $\to$ \en{Mean Shift} $\to$ зглобови}
Дакле, у овом систему Mean Shift првенствено обједињује класификоване пикселе у просторне предлоге, а не самостално одређује анатомску класу сваког пиксела.
\izvor{Прочитати §7.5.2, \en{Mean shift}, и §13.6.4, \en{Human body modeling and tracking}, из [S]; за тачан Kinect поступак [A1], §3.1--3.4, посебно \en{Joint position proposals}.}
\end{odgovor}

\section*{Литература и линкови}
\textbf{[S]} Richard Szeliski. \en{Computer Vision: Algorithms and Applications}. \textbf{2nd edition}. Springer, 2022. Ауторска страница са приступом електронском издању:
\url{https://szeliski.org/Book/}.

DOI: \url{https://doi.org/10.1007/978-3-030-34372-9}.

Ово је једина главна литература. Све референце на бројеве поглавља и одељака у скрипти односе се на \textbf{друго издање}. Прво издање из 2011. има другачији распоред поглавља. Странице ауторског PDF-а и штампаног издања могу се разликовати, па су испод одговора наведени одељци и њихови називи.

\medskip
\textbf{[A1]} Jamie Shotton et al. \en{Real-Time Human Pose Recognition in Parts from Single Depth Images}. CVPR, 2011.

Ауторски PDF: \url{https://www.microsoft.com/en-us/research/wp-content/uploads/2016/02/BodyPartRecognition.pdf}.

Овај додатни извор је употребљен за питање 50, јер раздваја класификацију делова тела од проналажења просторних модова за положаје зглобова. Најрелевантнији је §3.4, \en{Joint position proposals}; §3.1--3.3 објашњавају ознаке делова тела, дубинска обележја и шуме одлучивања.
\end{document}
